\documentclass[a4paper]{arrowhead}

\usepackage[yyyymmdd]{datetime}
\usepackage{etoolbox}
\usepackage[T1]{fontenc}
\usepackage[utf8]{inputenc}
\usepackage{multirow}
\usepackage{longtable}
\usepackage{amssymb}
\usepackage{float}
\usepackage{tikz}
\usetikzlibrary{positioning,arrows.meta,fit,calc,shapes.misc,shapes.geometric}

\renewcommand{\dateseparator}{-}
\setlength{\emergencystretch}{3em}
\lstset{basicstyle=\ttfamily\footnotesize,breaklines=true,columns=fullflexible,frame=single,rulecolor=\color{gray!50},backgroundcolor=\color{gray!6}}
\tikzset{
  ahsys/.style={draw=ArrowheadBlue,very thick,fill=ArrowheadBlue!6,rounded corners=2pt,minimum width=3.2cm,minimum height=1.1cm,align=center,font=\small\bfseries},
  ahext/.style={draw=ArrowheadGrey,thick,dashed,fill=gray!5,rounded corners=2pt,minimum width=3.2cm,minimum height=1.0cm,align=center,font=\small},
  ahsrv/.style={draw=ArrowheadPurple,thick,fill=ArrowheadPurple!8,rounded corners=6pt,minimum height=0.55cm,inner xsep=4pt,align=center,font=\scriptsize\ttfamily},
  ahcsm/.style={draw=ArrowheadTeal,thick,fill=ArrowheadTeal!8,rounded corners=6pt,minimum height=0.55cm,inner xsep=4pt,align=center,font=\scriptsize\ttfamily},
  ahdb/.style={draw=ArrowheadOrange,thick,fill=ArrowheadOrange!8,cylinder,shape border rotate=90,aspect=0.25,minimum height=0.9cm,align=center,font=\scriptsize},
  ahflow/.style={-{Stealth[length=2.2mm]},thick,ArrowheadGrey},
  ahlabel/.style={font=\scriptsize,fill=white,inner sep=1pt,align=center},
}

%% Special references
\newcommand{\mref}[1]{{\textcolor{ArrowheadPurple}{\hyperref[sec:model:#1]{#1}}}}
\newcommand{\oref}[1]{{\textcolor{ArrowheadBlue}{\hyperref[sec:operations:#1]{#1}}}}
\newcommand{\pdef}[1]{{\textcolor{ArrowheadGrey}{#1\label{sec:model:primitives:#1}\label{sec:model:primitives:#1s}\label{sec:model:primitives:#1es}}}}
\newcommand{\pref}[1]{{\textcolor{ArrowheadGrey}{\hyperref[sec:model:primitives:#1]{#1}}}}

\newrobustcmd\osubsection[3]{
  \addtocounter{subsection}{1}
  \addcontentsline{toc}{subsection}{\protect\numberline{\thesubsection}operation \textcolor{ArrowheadBlue}{#1}}
  \subsection*{
    \thesubsection\quad
    operation
    \textcolor{ArrowheadBlue}{#1}
    (\notblank{#2}{\mref{#2}}{})
    \notblank{#3}{: \mref{#3}}{}
  }
  \label{sec:operations:#1}
}
\newrobustcmd\msubsection[2]{
  \addtocounter{subsection}{1}
  \addcontentsline{toc}{subsection}{\protect\numberline{\thesubsection}#1 \textcolor{ArrowheadPurple}{#2}}
  \subsection*{\thesubsection\quad#1 \textcolor{ArrowheadPurple}{#2}}
  \label{sec:model:#2} \label{sec:model:#2s} \label{sec:model:#2es}
}

\begin{document}

%% Arrowhead Document Properties
\ArrowheadTitle{workflowWorkItem}
\ArrowheadServiceID{workflowWorkItem}
\ArrowheadType{Service Description}
\ArrowheadTypeShort{SD}
\ArrowheadVersion{5.4.0}
\ArrowheadDate{2026-11-23}
\ArrowheadAuthor{Documentation maintenance}
\ArrowheadStatus{REVIEW DRAFT}
\ArrowheadContact{Repository maintainers}
\ArrowheadFooter{\href{www.arrowhead.eu}{www.arrowhead.eu}}
\ArrowheadSetup
%%

%% Front Page
\begin{center}
  \vspace*{1cm}
  \huge{\arrowtitle}

  \vspace*{0.2cm}
  \LARGE{\arrowtype}
  \vspace*{1cm}

  \Large{Service ID: \textit{``\arrowid''}}
  \vspace*{\fill}

  

  \vspace*{1cm}
  \vspace*{\fill}

  % Front Page Abstract
  \begin{abstract}
    This document describes the \texttt{workflowWorkItem} service of the Eclipse Arrowhead workflow support systems, which lets executors claim work items and report their outcome. It is provided by the WorkflowManager system, service version 1.0.0, documentation edition 5.4.0.
  \end{abstract}

  \vspace*{1cm}
\end{center}

\newpage
%%

%% Table of Contents
\tableofcontents
\newpage
%%

\section{Overview}
\label{sec:overview}
The \texttt{workflowWorkItem} service is the interface between the WorkflowManager, which keeps the
books, and the WorkflowExecutors, which do the work. A \emph{work item} is one attempt to run one step of one session.
In pull mode executors \oref{claim} open work items they can run (the blackboard of~\cite{derhamy2018workflow}); in
every mode the executor holding a work item \oref{report}s what happened: the provider it was bound to, the result, or
the failure. The reports drive the transition rules S-Bind, S-Done, S-FailBind and S-FailInvoke.

The rest of this document is organized as follows.
In Section~\ref{sec:operations}, we describe the abstract message operations provided by the service.
In Section~\ref{sec:model}, we end the document by presenting the data types used by the mentioned operations.

\subsection{Significant Prior Art}
Derhamy's RetrieveOperation and StoreOperationResult interfaces~\cite{derhamy2018workflow}; the
Choreographer's \texttt{notify-step-done} and its result codes SUCCESS, ERROR, FATAL\_ERROR and ABORTED.

\subsection{How This Service Is Meant to Be Used}
A pull-mode executor calls \oref{claim} periodically with its capabilities and services. It receives the
first open work item it can run, or no content. After orchestration it reports \texttt{BOUND} (with the binding) or
\texttt{BIND\_FAILED}; after the invocation it first stores the result through \texttt{workflowContext.put-result} and
then reports \texttt{SUCCESS}, \texttt{ERROR} or \texttt{FATAL\_ERROR}. Push-mode executors receive their work items
through \texttt{workflowExecution.start} and report the same way.

\subsection{Important Delimitations and Dependencies}
\label{sec:delimitations}
\begin{itemize}
\item The claim is atomic: one executor per work item attempt. Only the holder (or a sysop) may report on a work item.
\item A claimed work item that is not reported within \texttt{workflow.pull.claim-timeout} (default 300\,s) fails the step with S-FailBind.
\item A report on a work item that is no longer active (already done, failed or aborted) is refused (409).
\end{itemize}

\paragraph{Authorization.} Each operation is authorized separately; executors need both. The holder check is made in addition.

\newpage
\section{Service Interface}
\label{sec:operations}
This section describes the interface of the \texttt{workflowWorkItem} service. Each subsection names an abstract operation,
an input type and an output type, in that order. The input type is named inside parentheses, while the output type is
preceded by a colon. Input and output types are only denoted when accepted or returned, respectively, by the operation
in question. All abstract data types named in this section are defined in Section~\ref{sec:model}. The concrete
interface (HTTP/TLS/JSON) is specified in the IDD of the same name.

\begin{figure}[ht!]
  \centering
  \resizebox{0.92\textwidth}{!}{
\begin{tikzpicture}
\node[ahsys] (prov) at (4.6,0) {WorkflowManager\\[-1pt]{\normalfont\scriptsize $\ll$provider$\gg$}};
\node[ahsrv] (srv) at (1.6,0) {workflowWorkItem};
\draw[thick,ArrowheadPurple] (srv.east) -- (prov.west);
\node[ahext] (c0) at (-3.0,0.00) {WorkflowExecutor\\[-1pt]{\scriptsize $\ll$consumer$\gg$}};
\draw[ahflow] (c0.east) -- (srv.west);
\end{tikzpicture}
  }
  \caption{SysML block view of the \texttt{workflowWorkItem} service: provider and typical consumers.}
  \label{fig:SD}
\end{figure}

The following interface operations are available.

{\small
\begin{longtable}{|p{4.4cm}|p{5.6cm}|p{5.8cm}|}

\hline
\rowcolor{gray!33} Operation & Input & Output \\ \hline
\endhead
\oref{claim} & \mref{ClaimRequest} & \mref{WorkItem} \\ \hline
\oref{report} & \mref{WorkItemReport} & -- \\ \hline
\end{longtable}
}

\osubsection{claim}{ClaimRequest}{WorkItem}
Atomically claims the first open work item the caller can run: its services cover $\theta(s)$ and its capabilities satisfy $\eta(s)$.
\osubsection{report}{WorkItemReport}{}
Reports on a held work item. \texttt{BOUND} completes S-Bind; \texttt{SUCCESS} fires S-Done with the stored result; \texttt{BIND\_FAILED} tries the next executor or fires S-FailBind; \texttt{ERROR} rebinds under a retry policy or fires S-FailInvoke; \texttt{FATAL\_ERROR} fires S-FailInvoke.

\newpage
\section{Information Model}
\label{sec:model}
Here, all data objects that can be part of the \texttt{workflowWorkItem} service are listed in alphabetic order. Each
subsection, which describes one type of object, begins with the \textit{struct} keyword (a collection of named fields,
each with its own data type), the \textit{enum} keyword (a closed set of values) or the \textit{alias} keyword. As a
complement to the explicitly defined types in this section, there is also a list of implicit primitive types in
Section~\ref{sec:model:primitives}.

\msubsection{struct}{ClaimRequest}
What the claiming executor can run. All fields are optional; the executor's identity comes from its certificate.

{\small
\begin{longtable}{|p{3.9cm}|p{4.4cm}|p{7.5cm}|}

\hline
\rowcolor{gray!33} Field & Type & Description \\ \hline
\endhead
\texttt{capabilities} & \pref{List}$<$\pref{String}$>$ & Declared capabilities $cap(w)$. Optional. \\ \hline
\texttt{services} & \pref{List}$<$\pref{ServiceDefinition}$>$ & Services the executor can run; empty or \texttt{*} means any. Optional. \\ \hline
\texttt{workstation} & \pref{String} & Workstation of the executor. Optional. \\ \hline
\end{longtable}
}

\msubsection{struct}{ExecutorRequirement}
$\eta(s)$: constraints on the executor that runs a step.

{\small
\begin{longtable}{|p{3.9cm}|p{4.4cm}|p{7.5cm}|}

\hline
\rowcolor{gray!33} Field & Type & Description \\ \hline
\endhead
\texttt{capability} & \pref{String} & A capability the executor must declare in \texttt{workflowCapabilities}. Optional. \\ \hline
\texttt{workstation} & \pref{String} & The executor's workstation. Optional. \\ \hline
\texttt{system} & \pref{SystemName} & A named executor system. Optional. \\ \hline
\end{longtable}
}

\msubsection{enum}{LeaseMode}
Role of a step for a shared resource.

{\small
\begin{longtable}{|p{3.9cm}|p{11.9cm}|}

\hline
\rowcolor{gray!33} Value & Meaning \\ \hline
\endhead
\texttt{reserve} & Reserves the resource; its result carries the lease token. \\ \hline
\texttt{use} & Uses the resource; the executor renews the lease immediately before the invocation. \\ \hline
\texttt{release} & Releases the resource. \\ \hline
\end{longtable}
}

\msubsection{struct}{LeasePolicy}
The role of a step for a shared resource.

{\small
\begin{longtable}{|p{3.9cm}|p{4.4cm}|p{7.5cm}|}

\hline
\rowcolor{gray!33} Field & Type & Description \\ \hline
\endhead
\texttt{resource} & \pref{Name} & A declared resource. \\ \hline
\texttt{mode} & \mref{LeaseMode} & \texttt{reserve}, \texttt{use} or \texttt{release}. \\ \hline
\end{longtable}
}

\msubsection{enum}{OnError}
Reaction to an \texttt{ERROR} report.

{\small
\begin{longtable}{|p{3.9cm}|p{11.9cm}|}

\hline
\rowcolor{gray!33} Value & Meaning \\ \hline
\endhead
\texttt{rebind} & Exclude the executor for this session and bind the step again. \\ \hline
\texttt{abort} & Fail the session. \\ \hline
\end{longtable}
}

\msubsection{struct}{Policy}
Execution policy of a step. Every field other than \texttt{lease} and \texttt{localCloudOnly} is an extension outside the proven model.

{\small
\begin{longtable}{|p{3.9cm}|p{4.4cm}|p{7.5cm}|}

\hline
\rowcolor{gray!33} Field & Type & Description \\ \hline
\endhead
\texttt{quantity} & \pref{Integer} & Number of executions; only 1 is accepted (check M01). Optional. \\ \hline
\texttt{timeoutSeconds} & \pref{Duration} & Extension x-timeout: the step fails when it runs longer. Optional. \\ \hline
\texttt{nominalDurationSeconds} & \pref{Duration} & Expected duration, used to check lease durations (check L04). Optional. \\ \hline
\texttt{retry} & \mref{Retry} & Extension x-rebind: behaviour on \texttt{ERROR}. Optional. \\ \hline
\texttt{lease} & \mref{LeasePolicy} & The step's role for a shared resource. Optional. \\ \hline
\texttt{localCloudOnly} & \pref{Boolean} & Whether orchestration may return providers of other clouds (default true). Optional. \\ \hline
\end{longtable}
}

\msubsection{enum}{ReportStatus}
Executor reports and the rules they cause.

{\small
\begin{longtable}{|p{3.9cm}|p{11.9cm}|}

\hline
\rowcolor{gray!33} Value & Meaning \\ \hline
\endhead
\texttt{BOUND} & Orchestration succeeded (completes S-Bind). \\ \hline
\texttt{BIND\_FAILED} & Orchestration returned no provider: the next executor is tried, otherwise S-FailBind. \\ \hline
\texttt{SUCCESS} & The provider answered (S-Done). \\ \hline
\texttt{ERROR} & Recoverable failure: x-rebind when the step has a retry policy, otherwise S-FailInvoke. \\ \hline
\texttt{FATAL\_ERROR} & Unrecoverable failure (S-FailInvoke). \\ \hline
\texttt{ABORTED} & The work item was aborted. \\ \hline
\end{longtable}
}

\msubsection{struct}{Retry}
Extension x-rebind.

{\small
\begin{longtable}{|p{3.9cm}|p{4.4cm}|p{7.5cm}|}

\hline
\rowcolor{gray!33} Field & Type & Description \\ \hline
\endhead
\texttt{maxAttempts} & \pref{Integer} & Maximum attempts of the step in one session. Optional. \\ \hline
\texttt{onError} & \mref{OnError} & What an \texttt{ERROR} report causes. Optional. \\ \hline
\end{longtable}
}

\msubsection{struct}{ServiceRequirement}
Requirements on the provider of the step's service, passed to orchestration.

{\small
\begin{longtable}{|p{3.9cm}|p{4.4cm}|p{7.5cm}|}

\hline
\rowcolor{gray!33} Field & Type & Description \\ \hline
\endhead
\texttt{serviceDefinition} & \pref{ServiceDefinition} & The service to consume. \\ \hline
\texttt{interfaces} & \pref{List}$<$\pref{String}$>$ & Accepted interface template names, e.g. \texttt{generic\_https}. Optional. \\ \hline
\texttt{minVersion} & \pref{String} & Lowest accepted service version. Optional. \\ \hline
\texttt{maxVersion} & \pref{String} & Highest accepted service version. Optional. \\ \hline
\texttt{metadata} & \pref{Map}$<$\pref{String}, \pref{JsonValue}$>$ & Metadata requirements on the service instance. Optional. \\ \hline
\texttt{qos} & \pref{Map}$<$\pref{String}, \pref{JsonValue}$>$ & QoS requirements forwarded to orchestration. Optional. \\ \hline
\texttt{requiresResults} & \pref{List}$<$\pref{Name}$>$ & Steps whose results the provider needs; each must be in $dep(s)$ (check D04). Optional. \\ \hline
\end{longtable}
}

\msubsection{struct}{WorkItem}
One attempt to run one step: what an executor needs to run it.

{\small
\begin{longtable}{|p{3.9cm}|p{4.4cm}|p{7.5cm}|}

\hline
\rowcolor{gray!33} Field & Type & Description \\ \hline
\endhead
\texttt{workItemId} & \pref{WorkItemId} & Identifier of the attempt. \\ \hline
\texttt{sessionId} & \pref{SessionId} & Session. \\ \hline
\texttt{recipe} & \pref{String} & \texttt{name@version} of the recipe. \\ \hline
\texttt{action} & \pref{Name} & Action of the step. \\ \hline
\texttt{step} & \pref{Name} & Step name. \\ \hline
\texttt{attempt} & \pref{Integer} & Attempt number from 1. \\ \hline
\texttt{service} & \mref{ServiceRequirement} & $\theta(s)$. \\ \hline
\texttt{executor} & \mref{ExecutorRequirement} & $\eta(s)$. Optional. \\ \hline
\texttt{params} & \pref{Map}$<$\pref{String}, \pref{JsonValue}$>$ & Static parameters of the step. Optional. \\ \hline
\texttt{input} & \pref{Template} & Input template $\kappa_s$. Optional. \\ \hline
\texttt{policy} & \mref{Policy} & Step policy. Optional. \\ \hline
\texttt{leaseDurationSeconds} & \pref{Duration} & Lease duration of the step's resource. Optional. \\ \hline
\texttt{parameters} & \pref{JsonObject} & Session parameters $\rho$ (included for convenience; also available from \texttt{workflowContext}). Optional. \\ \hline
\texttt{dependencyResults} & \pref{Map}$<$\pref{Name}, \pref{JsonValue}$>$ & $\xi|dep(s)$ (also available from \texttt{workflowContext}). Optional. \\ \hline
\end{longtable}
}

\msubsection{struct}{WorkItemReport}
A report on a held work item.

{\small
\begin{longtable}{|p{3.9cm}|p{4.4cm}|p{7.5cm}|}

\hline
\rowcolor{gray!33} Field & Type & Description \\ \hline
\endhead
\texttt{workItemId} & \pref{WorkItemId} & The work item. \\ \hline
\texttt{status} & \mref{ReportStatus} & What happened. \\ \hline
\texttt{binding} & \pref{JsonValue} & With \texttt{BOUND}: the provider the executor was given by orchestration. Optional. \\ \hline
\texttt{result} & \pref{JsonValue} & With \texttt{SUCCESS}: $\xi(s)$, the provider's response. Optional. \\ \hline
\texttt{message} & \pref{String} & Explanation, mandatory in practice for failures. Optional. \\ \hline
\end{longtable}
}


\subsection{Primitives}
\label{sec:model:primitives}
Types and structures mentioned throughout this document that are assumed to be available to implementations of this
service. The concrete interpretation of each of these types must be provided by any IDD document claiming to implement
this service.

{\small
\begin{longtable}{|p{3.2cm}|p{12.6cm}|}

\hline
\rowcolor{gray!33} Type & Description \\ \hline
\endhead
\pdef{Boolean} & One out of \texttt{true} or \texttt{false}. \\ \hline
\pdef{Duration} & A length of time in seconds. \\ \hline
\pdef{Integer} & A whole number. \\ \hline
\pdef{JsonObject} & A set of named values. \\ \hline
\pdef{JsonValue} & Any value: object, array, string, number, Boolean or null. \\ \hline
\pdef{List} & An ordered collection of items of the type given in angle brackets. \\ \hline
\pdef{Map} & A collection of key/value pairs of the types given in angle brackets. \\ \hline
\pdef{Name} & An identifier of a recipe, action, step, parameter or resource: a letter followed by at most 62 letters, digits, \texttt{-} or \texttt{\_}. \\ \hline
\pdef{ServiceDefinition} & The name of an Arrowhead service definition. Names are compared case-insensitively and without \texttt{-}, \texttt{\_} and spaces, so \texttt{temperature-read} and \texttt{temperatureRead} denote the same service. \\ \hline
\pdef{SessionId} & The identifier the WorkflowManager gives a session. \\ \hline
\pdef{String} & A chain of characters. \\ \hline
\pdef{SystemName} & The name of an Arrowhead system as it appears in its certificate and in the ServiceRegistry. \\ \hline
\pdef{Template} & A value in which strings of the form \texttt{\$\{param.$p$\}} or \texttt{\$\{result.$s$.path\}} are replaced by a session parameter or by (a part of) a stored step result. \\ \hline
\pdef{WorkItemId} & The identifier of one attempt to run one step of one session. \\ \hline
\end{longtable}
}


\newpage
\bibliographystyle{IEEEtran}
\bibliography{bibliography}

\newpage

\section{Revision History}
\subsection{Amendments}

\noindent\begin{tabularx}{\textwidth}{| p{1cm} | p{2.2cm} | p{1.6cm} | X | p{3.6cm} |} \hline
\rowcolor{gray!33} No. & Date & Version & Subject of Amendments & Author \\ \hline
1 & 2026-11-23 & \arrowversion & Initial edition: service description of workflowWorkItem, written from the Java reference implementation (ah-Workflow-java). & Documentation maintenance \\ \hline
\end{tabularx}

\subsection{Quality Assurance}

\noindent\begin{tabularx}{\textwidth}{| p{1cm} | p{2.2cm} | p{1.6cm} | X |} \hline
\rowcolor{gray!33} No. & Date & Version & Approved by \\ \hline
1 & & \arrowversion & Pending review \\ \hline
\end{tabularx}

\end{document}
